\section{Skeleton-averaged representation}
\label{sec:cprime}

We now average all continuous variables inside a finite decorated skeleton before taking absolute values. This produces a countable signed series whose absolute convergence can be checked directly in the parabolic H\"older space.

Let
\begin{equation*}
  M=\norm L_X,
  \qquad
  a=\abs\lambda C_{\cD}(\alpha,T)M.
\end{equation*}
For $\tau\in\cT$, define $F_\tau$ by~\eqref{eq:tree-profile-recursion}. If $\cT_n$ denotes the trees with $n$ internal vertices, then
\begin{equation*}
  \abs{\cT_n}
  =
  C_n
  =
  \frac1{n+1}\binom{2n}{n}.
\end{equation*}
The bilinear bound~\eqref{eq:D-bilinear-bound} gives
\begin{equation}
  \norm{F_\tau}_X
  \leq
  Ma^n,
  \qquad \tau\in\cT_n.
  \label{eq:Cprime-tree-bound}
\end{equation}

\begin{theorem}[Skeleton-averaged $L^1$ representation, Theorem C-prime]
\label{thm:cprime}
Fix $0<\alpha<1$, $T>0$, $\lambda\in\R$, and $\phi\in C^{2+\alpha}(\Torus)$. Let
\begin{equation*}
  X=X_{\alpha,T},
  \qquad
  M=\norm{P_\cdot\phi''}_{X},
\end{equation*}
and let $C_{\cD}(\alpha,T)$ be a valid constant in the bilinear estimate~\eqref{eq:D-bilinear-bound}. Put
\begin{equation*}
  a
  =
  \abs\lambda C_{\cD}(\alpha,T)M.
\end{equation*}
Assume
\begin{equation}
  4a<1.
  \label{eq:Cprime-smallness}
\end{equation}
Then
\begin{equation}
  z_*
  =
  \sum_{\tau\in\cT}F_\tau
  \label{eq:Cprime-series}
\end{equation}
converges absolutely in $X$. It satisfies
\begin{equation*}
  z_*
  =
  P_\cdot\phi''
  +\lambda\cD(z_*^2).
\end{equation*}
Define
\begin{equation*}
  R_*(a)
  =
  M\frac{1-\sqrt{1-4a}}{2a},
  \qquad
  R_*(0)=M.
\end{equation*}
Then $z_*$ is the unique fixed point of $z\mapsto P_\cdot\phi''+\lambda\cD(z^2)$ in the closed $X$-ball of radius $R_*(a)$.

Let $\pi$ be any probability mass function on $\cT$ with $\pi(\tau)>0$ for every finite tree. If $S\sim\pi$ and
\begin{equation}
  \widehat z(t,x)
  =
  \frac{F_S(t,x)}{\pi(S)},
  \label{eq:Cprime-estimator}
\end{equation}
then, for every $(t,x)\in[0,T]\times\Torus$,
\begin{equation*}
  \widehat z(t,x)\in L^1,
  \qquad
  \E[\widehat z(t,x)]=z_*(t,x).
\end{equation*}
Moreover, defining
\begin{equation*}
  v_*(t)
  =
  P_t\phi
  +\lambda\int_0^tP_{t-s}[z_*(s)^2]\,ds,
\end{equation*}
one has
\begin{equation*}
  v_*\in C^{1+\alpha/2,2+\alpha}([0,T]\times\Torus),
  \qquad
  v_{*,xx}=z_*,
\end{equation*}
so $v_*$ is a classical solution of~\eqref{eq:hessian-forward}.
\end{theorem}

A data-only sufficient condition follows from the heat-lift estimate
\begin{equation*}
  \norm{P_\cdot\phi''}_X
  \leq
  (1+\E\abs Z^\alpha)\norm{\phi''}_{C^\alpha}.
\end{equation*}

The theorem is not defined by conditioning an unresolved infinite raw patch functional. Each $F_\tau$ is first defined by the deterministic Duhamel recursion. At finite cutoff it can be recovered by averaging the continuous proposal variables inside the fixed skeleton, but the infinite theorem starts from the absolutely summable deterministic family.

\begin{proof}
The estimate~\eqref{eq:Cprime-tree-bound} follows by induction on the number of internal vertices. It is immediate for the leaf. If $\tau=[\tau_1,\tau_2]$ and the two subtrees have $n_1,n_2$ internal vertices, then $n=n_1+n_2+1$ and
\begin{align*}
  \norm{F_\tau}_X
  &\leq
  \abs\lambda C_{\cD}
  \norm{F_{\tau_1}}_X
  \norm{F_{\tau_2}}_X \\
  &\leq
  \abs\lambda C_{\cD}
  M^2a^{n_1+n_2}
  =
  Ma^n.
\end{align*}
Therefore
\begin{equation}
  \sum_{\tau\in\cT}\norm{F_\tau}_X
  \leq
  M\sum_{n=0}^\infty C_na^n.
  \label{eq:Cprime-Catalan-sum}
\end{equation}
Under~\eqref{eq:Cprime-smallness}, $a<1/4$, and the Catalan generating function gives
\begin{equation}
  \sum_{n=0}^\infty C_na^n
  =
  \frac{1-\sqrt{1-4a}}{2a}.
  \label{eq:Catalan-generating}
\end{equation}
Thus~\eqref{eq:Cprime-series} converges absolutely in $X$.

The non-leaf planar trees are in bijection with ordered pairs $(\tau_1,\tau_2)$ by removing the root. Since the series in~\eqref{eq:Cprime-Catalan-sum} is absolutely convergent and $(f,g)\mapsto\lambda\cD(fg)$ is a continuous bilinear map on $X$, the corresponding double series is absolutely convergent. Hence
\begin{align*}
  z_*-L
  &=
  \sum_{\tau_1,\tau_2}
  \lambda\cD(F_{\tau_1}F_{\tau_2}) \\
  &=
  \lambda\cD\br{
    \pa{\sum_{\tau_1}F_{\tau_1}}
    \pa{\sum_{\tau_2}F_{\tau_2}}
  }
  =
  \lambda\cD(z_*^2).
\end{align*}
The absolute convergence is load-bearing here: it justifies the Cauchy product and the exchange of the two tree sums.

The Catalan identity gives
\begin{equation}
  R_*(a)
  =
  M+\abs\lambda C_{\cD}R_*(a)^2.
  \label{eq:Cprime-radius-identity}
\end{equation}
For
\begin{equation*}
  \Phi(z)=L+\lambda\cD(z^2),
\end{equation*}
~\eqref{eq:Cprime-radius-identity} shows that $\Phi$ maps the closed $X$-ball of radius $R_*(a)$ into itself. If $\norm z_X,\norm w_X\leq R_*(a)$, then
\begin{equation*}
  \norm{\Phi(z)-\Phi(w)}_X
  \leq
  2\abs\lambda C_{\cD}R_*(a)\norm{z-w}_X.
\end{equation*}
Using the explicit value of $R_*(a)$,
\begin{equation*}
  2\abs\lambda C_{\cD}R_*(a)
  =
  1-\sqrt{1-4a}<1.
\end{equation*}
Thus $z_*$ is the unique fixed point in this ball.

Define $v_*$ as in the theorem statement. Differentiating twice in space and using the fixed-point equation gives $v_{*,xx}=z_*$. Standard parabolic Schauder regularity yields the claimed regularity and therefore the classical solution property.

It remains to prove the probabilistic statement. Sample $S\sim\pi$ and define~\eqref{eq:Cprime-estimator}. Then
\begin{align*}
  \E\abs{\widehat z(t,x)}
  &=
  \sum_{\tau\in\cT}\abs{F_\tau(t,x)} \\
  &\leq
  \sum_{\tau\in\cT}\norm{F_\tau}_X
  <\infty.
\end{align*}
The expectation is therefore absolutely convergent, and the same summability gives
\begin{equation*}
  \E[\widehat z(t,x)]
  =
  \sum_{\tau\in\cT}F_\tau(t,x)
  =
  z_*(t,x).
\end{equation*}
The skeleton proposal cancels from the first moment exactly as in ordinary importance sampling, although it still affects higher moments and computational efficiency.
\end{proof}

\begin{corollary}[Overlap with the deterministic solution]
\label{cor:Cprime-deterministic-overlap}
Fix $0<\alpha<1$, $T>0$, $\lambda\in\R$, and $\phi\in C^{2+\alpha}(\Torus)$. Assume the C-prime condition~\eqref{eq:Cprime-smallness} and the deterministic smallness condition
\begin{equation*}
  \abs\lambda C_{\mathrm{Sch}}(\alpha,T)
  \norm\phi_{C^{2+\alpha}}
  \leq\frac18.
\end{equation*}
Let $z_*$ be the C-prime fixed point, and assume in addition that
\begin{equation*}
  \abs{\lambda z_*}\leq\frac18
  \qquad\text{on }[0,T]\times\Torus.
\end{equation*}
Then the C-prime classical solution coincides with the solution produced by~\cref{thm:deterministic-iteration}.
\end{corollary}

\begin{proof}
Both constructions give classical solutions of~\eqref{eq:hessian-forward}. The displayed bound puts the C-prime solution in the uniqueness class of~\cref{thm:deterministic-iteration}, so the two solutions coincide.
\end{proof}
